\documentclass[aspectratio=169,10pt]{beamer}
\makeatletter\def\input@path{{../estilo/}{estilo/}}\makeatother
\usepackage{devmobile}
\modulo{14}{Tópicos Avançados e Projeto Final}{IA no dispositivo e na nuvem, segurança, monetização, outras telas, carreira e o projeto integrador}

\begin{document}
\capa

\begin{frame}{Objetivos do módulo}
  \begin{conceito}[Ao final deste módulo você será capaz de]
    \begin{itemize}
      \item Integrar modelos de aprendizado de máquina no app, no aparelho e na nuvem.
      \item Aplicar as principais práticas de segurança mobile (OWASP MASVS).
      \item Conhecer os modelos de monetização e as APIs da Google Play.
      \item Levar o app a widgets, relógios e carros.
      \item Planejar, executar e apresentar o projeto final do curso.
    \end{itemize}
  \end{conceito}
\end{frame}

\sumario

% =============================================================================
\section{Inteligência artificial no app}

\begin{frame}{Onde roda o modelo?}
  \begin{columns}[c]
    \begin{column}{0.54\textwidth}
      \centering
      \begin{tikzpicture}[node distance=4mm]
        \node (cel) {\celular[0.5]{\tiny modelo\\ local}};
        \node[caixaambar,right=12mm of cel,yshift=8mm,text width=3.4cm] (dev) {\textbf{No aparelho}\\ \tiny LiteRT, ML Kit, Gemini Nano};
        \node[caixaescura,right=12mm of cel,yshift=-8mm,text width=3.4cm] (nuv) {\textbf{Na nuvem}\\ \tiny Firebase AI Logic, APIs de LLM};
        \draw[seta] (cel.east) -- (dev.west);
        \draw[seta] (cel.east) -- (nuv.west);
      \end{tikzpicture}
    \end{column}
    \begin{column}{0.44\textwidth}
      \scriptsize
      \begin{tabularx}{\linewidth}{@{}l L L@{}}
        \toprule
        & \cabecalho{Aparelho} & \cabecalho{Nuvem} \\
        \midrule
        Latência & muito baixa & depende da rede \\
        Offline & sim & não \\
        Privacidade & dados ficam no aparelho & dados saem \\
        Custo & nenhum por uso & por requisição \\
        Tamanho do modelo & pequeno (KB--MB) & grande \\
        \bottomrule
      \end{tabularx}
    \end{column}
  \end{columns}
  \vspace{4pt}
  \begin{columns}[T]
    \begin{column}{0.32\textwidth}
      \begin{conceito}[ML Kit]
        Modelos prontos do Google: texto em imagem (OCR), códigos de barras, rostos, rótulos, tradução.
      \end{conceito}
    \end{column}
    \begin{column}{0.32\textwidth}
      \begin{conceito}[LiteRT]
        Novo nome do \textbf{TensorFlow Lite}: roda o \textbf{seu} modelo (\arq{.tflite}) com aceleração por GPU/NPU.
      \end{conceito}
    \end{column}
    \begin{column}{0.32\textwidth}
      \begin{conceito}[Gemini Nano]
        LLM do sistema em aparelhos compatíveis, acessível pelas APIs de IA generativa do ML Kit.
      \end{conceito}
    \end{column}
  \end{columns}
\end{frame}

\begin{frame}[fragile]{Detectando anomalias nas leituras}
  \begin{columns}[T]
    \begin{column}{0.5\textwidth}
\begin{codigo}[title=1) Linha de base estatística (z-score)]{Kotlin}
fun Flow<Double>.anomalias(
  janela: Int = 60, limiar: Double = 3.0
): Flow<Double> = flow {
  val buf = ArrayDeque<Double>()
  collect { x ->
    if (buf.size == janela) {
      val m = buf.average()
      val dp = sqrt(buf.sumOf { (it - m).pow(2) } / janela)
      if (dp > 0 && abs(x - m) / dp > limiar) emit(x)
      buf.removeFirst()
    }
    buf.addLast(x)
  }
}
\end{codigo}
      \begin{dica}
        Sempre compare o modelo de IA com uma linha de base simples como esta. Às vezes ela basta.
      \end{dica}
    \end{column}
    \begin{column}{0.48\textwidth}
\begin{codigo}[title=2) Modelo próprio com LiteRT]{Kotlin}
class DetectorAnomalia(ctx: Context) {
  private val modelo = Interpreter(
    FileUtil.loadMappedFile(ctx, "anomalia.tflite"))

  // janela de 60 leituras normalizadas -> prob.
  fun probabilidade(janela: FloatArray): Float {
    val entrada = arrayOf(janela)          // [1, 60]
    val saida = Array(1) { FloatArray(1) } // [1, 1]
    modelo.run(entrada, saida)
    return saida[0][0]
  }
}
\end{codigo}
\begin{terminal}[title=Treino em Python (Keras) e conversão]
conv = tf.lite.TFLiteConverter.from_keras_model(m)
conv.optimizations = [tf.lite.Optimize.DEFAULT]
open("anomalia.tflite", "wb").write(conv.convert())
\end{terminal}
    \end{column}
  \end{columns}
\end{frame}

\begin{frame}{IA na nuvem e boas práticas}
  \begin{columns}[T]
    \begin{column}{0.5\textwidth}
      \begin{conceito}[Firebase AI Logic]
        SDK para chamar modelos Gemini a partir do app \textbf{sem expor chaves}: o Firebase intermedeia a chamada e o \textbf{App Check} garante que só o seu app usa a cota.
      \end{conceito}
      \begin{exemplo}[Ideias para o LabMonitor]
        \begin{itemize}
          \item ``Explique este alerta'' a partir das últimas leituras.
          \item Fotografar a placa do equipamento e preencher o cadastro (OCR).
          \item Resumo semanal de uso das bancadas em linguagem natural.
        \end{itemize}
      \end{exemplo}
    \end{column}
    \begin{column}{0.48\textwidth}
      \begin{atencao}[Cuidados]
        \begin{itemize}
          \item Respostas de LLM podem estar erradas: mostre como sugestão e deixe o usuário confirmar.
          \item Não envie dados pessoais sem necessidade e sem informar o usuário (LGPD).
          \item Defina limites de uso e trate falhas e lentidão da rede.
          \item Quantize modelos locais (\emph{int8}) para reduzir tamanho e consumo.
        \end{itemize}
      \end{atencao}
      \begin{dica}
        Valide o desempenho do modelo no aparelho mais fraco que seu público usa, não só no seu celular.
      \end{dica}
    \end{column}
  \end{columns}
\end{frame}

% =============================================================================
\section{Segurança mobile}

\begin{frame}{OWASP MASVS: os pilares}
  \footnotesize
  \begin{tabularx}{\textwidth}{@{}l L@{}}
    \toprule
    \cabecalho{Área} & \cabecalho{Práticas no Android} \\
    \midrule
    \textbf{Armazenamento} & dados sensíveis só no armazenamento interno; chaves no \textbf{Android Keystore}; nada sensível em logs ou backups \\
    \textbf{Criptografia} & algoritmos padrão (AES-GCM, RSA/EC) via Keystore; nunca criptografia ``caseira'' \\
    \textbf{Autenticação} & tokens de curta duração; biometria com \cod{BiometricPrompt}; \emph{passkeys} com Credential Manager \\
    \textbf{Rede} & HTTPS sempre; \emph{network security config}; \emph{pinning} só com plano de rotação \\
    \textbf{Plataforma} & expor o mínimo de componentes; validar \emph{intents} e \emph{deep links}; WebView sem JavaScript desnecessário \\
    \textbf{Código} & dependências atualizadas; R8; nenhum segredo no APK \\
    \textbf{Resiliência} & \textbf{Play Integrity API} para verificar app e aparelho no servidor \\
    \bottomrule
  \end{tabularx}
  \vspace{4pt}
  \begin{destaque}
    O cliente nunca é confiável: toda regra importante precisa ser verificada também no servidor.
  \end{destaque}
\end{frame}

% =============================================================================
\section{Monetização e outras telas}

\begin{frame}{Modelos de monetização}
  \begin{columns}[T]
    \begin{column}{0.5\textwidth}
      \scriptsize
      \begin{tabularx}{\linewidth}{@{}l L@{}}
        \toprule
        \cabecalho{Modelo} & \cabecalho{Como funciona} \\
        \midrule
        Pago & usuário paga para baixar \\
        Freemium & grátis com recursos pagos (compras no app) \\
        Assinatura & cobrança recorrente por acesso \\
        Anúncios & receita por exibição (AdMob) \\
        B2B / institucional & a organização paga pelo app e pelo serviço \\
        Gratuito & projetos acadêmicos, públicos e de extensão \\
        \bottomrule
      \end{tabularx}
    \end{column}
    \begin{column}{0.48\textwidth}
      \begin{conceito}[Google Play Billing]
        Bens digitais vendidos dentro de apps distribuídos pela Play usam o sistema de faturamento da Play, que retém uma taxa de serviço. Bens físicos e serviços prestados fora do app seguem outras regras.
      \end{conceito}
      \begin{atencao}
        Valide as compras \textbf{no servidor} com a API da Play antes de liberar o conteúdo.
      \end{atencao}
      \begin{dica}
        Para apps institucionais, o modelo comum é gratuito, sem anúncios, com o custo coberto pelo projeto.
      \end{dica}
    \end{column}
  \end{columns}
\end{frame}

\begin{frame}{Além do celular}
  \begin{columns}[T]
    \begin{column}{0.32\textwidth}
      \centering
      \begin{tikzpicture}
        \draw[petroleo,rounded corners=4pt,fill=petroleo!6] (0,0) rectangle (2.6,1.4);
        \node[font=\tiny\bfseries,text=petroleo] at (1.3,1.15) {LabMonitor};
        \node[font=\large\bfseries,text=ambar] at (1.3,0.65) {24,8\,°C};
        \node[font=\tiny,text=cinzamedio] at (1.3,0.25) {Bancada 3 $\cdot$ há 1 min};
      \end{tikzpicture}
      \begin{conceito}[Widgets (Glance)]
        Informação na tela inicial, escrita com uma API no estilo Compose.
      \end{conceito}
    \end{column}
    \begin{column}{0.32\textwidth}
      \centering
      \begin{tikzpicture}
        \draw[petroleo,line width=1.4pt,fill=petroleoesc] (0,0) circle (0.72);
        \node[font=\small\bfseries,text=white] at (0,0.1) {24,8°};
        \node[font=\tiny,text=ambarclaro] at (0,-0.3) {B3};
      \end{tikzpicture}
      \begin{conceito}[Wear OS]
        Apps e \emph{tiles} no relógio com Compose for Wear OS; ideal para alertas rápidos.
      \end{conceito}
    \end{column}
    \begin{column}{0.32\textwidth}
      \centering
      \begin{tikzpicture}
        \draw[petroleo,rounded corners=3pt,fill=petroleo!6] (0,0) rectangle (2.8,1.4);
        \fill[petroleo] (0,1.1) rectangle (2.8,1.4);
        \node[font=\tiny,white] at (1.4,1.25) {Android Auto};
        \node[font=\tiny,text=cinzatexto] at (1.4,0.6) {lista simples, voz};
      \end{tikzpicture}
      \begin{conceito}[Carro e TV]
        Android Auto, Automotive e TV usam \emph{templates} e regras de segurança próprias.
      \end{conceito}
    \end{column}
  \end{columns}
  \begin{dica}
    A arquitetura em camadas permite reaproveitar repositórios e casos de uso em todas essas superfícies; só a camada de UI muda.
  \end{dica}
\end{frame}

% =============================================================================
\section{Carreira e portfólio}

\begin{frame}{Do curso ao mercado}
  \begin{columns}[T]
    \begin{column}{0.48\textwidth}
      \begin{conceito}[Portfólio que convence]
        \begin{itemize}
          \item 2--3 apps completos no GitHub, com README, capturas e arquitetura explicada.
          \item Pelo menos um app publicado na Play (mesmo em teste aberto).
          \item Testes e CI visíveis no repositório.
          \item Um projeto com hardware (BLE/IoT) é um diferencial raro.
        \end{itemize}
      \end{conceito}
    \end{column}
    \begin{column}{0.48\textwidth}
      \begin{dica}[Continue aprendendo]
        \begin{itemize}
          \item Código-fonte do \emph{Now in Android} e dos \emph{samples} oficiais.
          \item Blog \emph{Android Developers} e as notas de versão do Jetpack.
          \item Comunidades: GDG, Kotlin User Groups, eventos locais.
          \item Contribua com bibliotecas de código aberto.
        \end{itemize}
      \end{dica}
      \begin{atencao}[Assistentes de IA]
        Use-os para acelerar, mas entenda e revise cada linha: você responde pelo código que publica.
      \end{atencao}
    \end{column}
  \end{columns}
\end{frame}

% =============================================================================
\section{Projeto final}

\begin{frame}{Projeto final: requisitos}
  \begin{columns}[T]
    \begin{column}{0.5\textwidth}
      \begin{conceito}[Requisitos mínimos]
        \begin{enumerate}
          \item Pelo menos 4 telas com Navigation 3 e Material 3.
          \item Arquitetura MVVM com repositórios e Hilt.
          \item Persistência local (Room e DataStore).
          \item Integração com API REST ou Firebase.
          \item Funcionamento offline nas telas principais.
          \item Pelo menos um recurso do dispositivo (BLE, câmera, sensores, localização).
          \item Testes unitários dos ViewModels e casos de uso.
          \item Acessibilidade verificada e tema escuro.
        \end{enumerate}
      \end{conceito}
    \end{column}
    \begin{column}{0.48\textwidth}
      \begin{projeto}[Sugestões de tema]
        \begin{itemize}
          \item Concluir e publicar o \textbf{LabMonitor}.
          \item Controle de empréstimo de equipamentos do laboratório via QR code.
          \item Medidor de consumo de energia com ESP32 e sensor de corrente.
          \item App de acompanhamento de estágio ou TCC.
          \item Coletor de dados de vibração de motores pelo acelerômetro.
        \end{itemize}
      \end{projeto}
      \begin{dica}
        Resolva um problema real de alguém que você conhece.
      \end{dica}
    \end{column}
  \end{columns}
\end{frame}

\begin{frame}{Entregáveis, cronograma e avaliação}
  \begin{columns}[T]
    \begin{column}{0.48\textwidth}
      \scriptsize
      \begin{tabularx}{\linewidth}{@{}l L@{}}
        \toprule
        \cabecalho{Semana} & \cabecalho{Marco} \\
        \midrule
        1 & proposta: problema, público, telas (protótipo) e arquitetura \\
        2 & navegação, UI e dados locais \\
        3 & rede/Firebase, recurso do dispositivo e offline \\
        4 & testes, acessibilidade, release e apresentação \\
        \bottomrule
      \end{tabularx}
      \begin{conceito}[Entregáveis]
        Repositório GitHub com README, APK/AAB de release, vídeo de 3 minutos e apresentação de 10 minutos.
      \end{conceito}
    \end{column}
    \begin{column}{0.5\textwidth}
      \scriptsize
      \begin{tabularx}{\linewidth}{@{}L r@{}}
        \toprule
        \cabecalho{Critério de avaliação} & \cabecalho{Peso} \\
        \midrule
        Funcionalidades e requisitos atendidos & 25\,\% \\
        Arquitetura e qualidade do código & 20\,\% \\
        Experiência do usuário, UI e acessibilidade & 15\,\% \\
        Testes automatizados e CI & 15\,\% \\
        Uso de recursos do dispositivo / integração & 10\,\% \\
        Documentação (README, diagramas) & 5\,\% \\
        Apresentação e demonstração & 10\,\% \\
        \midrule
        \textbf{Total} & \textbf{100\,\%} \\
        \bottomrule
      \end{tabularx}
      \begin{dica}
        Pesos sugeridos: ajuste conforme o plano de ensino da disciplina.
      \end{dica}
    \end{column}
  \end{columns}
\end{frame}

% =============================================================================
\section{Fechamento do curso}

\begin{frame}{A jornada completa}
  \begin{center}
  \begin{tikzpicture}[x=1.08cm,
      etapa/.style={caixa,minimum width=1.75cm,minimum height=9mm,font=\tiny,text width=1.6cm}]
    \foreach \i/\t/\s in {0/Ecossistema/caixa,1/Ambiente/caixa,2/Kotlin/caixa,3/Coroutines/caixa,4/Compose/caixaambar,
                          5/Estado e navegação/caixaambar,6/Arquitetura/caixaambar}{
      \node[etapa,\s] (a\i) at (\i*1.9,0) {M\i\\ \t};}
    \foreach \i/\t/\s [count=\k from 7] in {0/Persistência/caixapetroleo,1/Rede e APIs/caixapetroleo,2/Firebase/caixapetroleo,
                          3/Dispositivo e IoT/caixapetroleo,4/Testes/caixaescura,5/Publicação/caixaescura,6/Multiplataforma/caixaescura}{
      \node[etapa,\s] (b\i) at (\i*1.9,-1.5) {M\k\\ \t};}
    \foreach \i [evaluate=\i as \j using int(\i+1)] in {0,...,5}{\draw[seta] (a\i) -- (a\j); \draw[seta] (b\i) -- (b\j);}
    \draw[seta,rounded corners] (a6.south) -- ++(0,-2mm) -| (b0.north);
    \node[caixaambarforte,minimum width=12cm,minimum height=8mm] at (5.7,-2.9) {\textbf{M14 --- IA, segurança e projeto final: o LabMonitor publicado}};
  \end{tikzpicture}
  \end{center}
  \begin{resumo}[O que você construiu]
    Um app Android completo e publicado, com arquitetura profissional, dados locais e remotos, integração com hardware, testes e entrega contínua, além da base para levá-lo a outras plataformas.
  \end{resumo}
\end{frame}

\begin{frame}{Exercícios finais}
  \begin{exercicio}[Exercício 1 --- IA]
    Gere 24\,h de leituras simuladas com picos artificiais e compare o detector por z-score com um modelo treinado em Python e convertido para LiteRT: quantos picos cada um encontra?
  \end{exercicio}
  \begin{exercicio}[Exercício 2 --- Segurança]
    Faça uma revisão do LabMonitor com o checklist MASVS do slide e registre ao menos três correções em \emph{issues} no GitHub.
  \end{exercicio}
  \begin{exercicio}[Exercício 3 --- Widget]
    Crie um widget Glance que mostre a última temperatura da bancada favorita e abra o app ao ser tocado.
  \end{exercicio}
  \begin{projeto}[Projeto final]
    Escolha o tema, entregue a proposta na semana 1 e siga o cronograma de 4 semanas.
  \end{projeto}
\end{frame}

\begin{frame}{Para aprofundar}
  \small
  \begin{itemize}
    \item IA no Android: \url{https://developer.android.com/ai}
    \item LiteRT: \url{https://ai.google.dev/edge/litert} \quad ML Kit: \url{https://developers.google.com/ml-kit}
    \item Firebase AI Logic: \url{https://firebase.google.com/docs/ai-logic}
    \item OWASP MASVS: \url{https://mas.owasp.org/MASVS/}
    \item Segurança no Android: \url{https://developer.android.com/privacy-and-security/security-tips}
    \item Google Play Billing: \url{https://developer.android.com/google/play/billing}
    \item Glance (widgets): \url{https://developer.android.com/develop/ui/compose/glance}
  \end{itemize}
\end{frame}

\begin{frame}[plain,noframenumbering]
  \begin{tikzpicture}[remember picture,overlay]
    \fill[petroleo] (current page.south west) rectangle (current page.north east);
    \fill[ambar] ([yshift=-0.9cm]current page.west) rectangle ([yshift=-0.82cm]current page.east);
    \node[font=\Huge\bfseries,text=white] at ([yshift=0.7cm]current page.center) {Obrigado!};
    \node[font=\large,text=ambarclaro] at ([yshift=-0.2cm]current page.center) {Fim do curso \nomecurso};
    \node[font=\small,text=white!80,text width=11cm,align=center] at ([yshift=-1.7cm]current page.center)
      {\autorcurso\\ \instituicaocurso\\ \departamentocurso};
  \end{tikzpicture}
\end{frame}

\end{document}
